\documentclass{article}

\usepackage{amsmath,amssymb}
\usepackage{xurl}
\usepackage[hidelinks]{hyperref}
\setlength{\emergencystretch}{2em}

% Shared commands for notes in docs/paper-gaps/.

\DeclareUnicodeCharacter{2080}{\ifmmode{}_{0}\else\textsubscript{0}\fi}
\DeclareUnicodeCharacter{2081}{\ifmmode{}_{1}\else\textsubscript{1}\fi}
\DeclareUnicodeCharacter{2082}{\ifmmode{}_{2}\else\textsubscript{2}\fi}
\DeclareUnicodeCharacter{2099}{\ifmmode{}_{n}\else\textsubscript{n}\fi}

\newcommand{\Lean}{\textsc{Lean}}
\ExplSyntaxOn
\NewDocumentCommand{\leanid}{m}{
  \ifmmode
    \textsf{\detokenize{#1}}
  \else
    \group_begin:
    \urlstyle{sf}
    \tl_set:Nn \l_tmpa_tl {#1}
    \tl_replace_all:Nnn \l_tmpa_tl {\_} {_}
    \exp_args:NV \path \l_tmpa_tl
    \group_end:
  \fi
}
\ExplSyntaxOff
\providecommand{\tr}{\operatorname{tr}}
\newcommand{\tnleanIssueBase}{https://github.com/LionSR/TNLean/issues/}
\newcommand{\tnleanPrBase}{https://github.com/LionSR/TNLean/pull/}
\newcommand{\tnleanDocBase}{https://sirui-lu.com/TNLean/docs/}
\newcommand{\ghissue}[1]{\href{\tnleanIssueBase#1}{GitHub issue \##1}}
\newcommand{\ghpr}[1]{\href{\tnleanPrBase#1}{PR \##1}}
\newcommand{\doclink}[1]{\href{\tnleanDocBase#1.html}{\path{#1}}}

% Dirac notation and the identity operator, as in blueprint/src/macros/common.tex.
% \providecommand keeps note-local definitions authoritative.
\usepackage{dsfont}
\providecommand{\ket}[1]{|#1\rangle}
\providecommand{\bra}[1]{\langle #1|}
\providecommand{\braket}[2]{\langle #1 | #2 \rangle}
\providecommand{\ketbra}[2]{|#1\rangle\!\langle #2|}
\providecommand{\Id}{\mathds{1}}
\providecommand{\one}{\mathds{1}}
\providecommand{\C}{\mathbb{C}}
\providecommand{\MN}[1]{M_{#1}(\C)}
\providecommand{\MD}{M_D(\C)}

% External referencing for paper sources.  The build helper writes sanitized
% auxiliary files under build/paper-aux/ and excludes duplicated source labels.
\usepackage{xr}

% Import a generated paper auxiliary file with a caller-supplied label prefix.
% The candidate paths support builds from docs/paper-gaps/, the repository root,
% and build/paper-aux/ itself.
\newcommand{\importpaperaux}[2]{%
  \IfFileExists{../../build/paper-aux/#2.aux}{%
    \externaldocument[#1]{../../build/paper-aux/#2}%
  }{%
    \IfFileExists{build/paper-aux/#2.aux}{%
      \externaldocument[#1]{build/paper-aux/#2}%
    }{%
      \IfFileExists{#2.aux}{%
        \externaldocument[#1]{#2}%
      }{}%
    }%
  }%
}

% General external reference macros
\newcommand{\paperref}[2]{\ref{#1-#2}}
\newcommand{\papereqref}[2]{(\ref{#1-#2})}

% CPSV16 source (1606.00608 / MPDO-22-12-17-2.tex) external document and shortcuts
\importpaperaux{cpsv16-}{1606.00608/MPDO-22-12-17-2-tnlean}
\newcommand{\cpsvref}[1]{\paperref{cpsv16}{#1}}
\newcommand{\cpsveqref}[1]{\papereqref{cpsv16}{#1}}

% Verdict marker: \gapnote{<kind>}{<status>}. Kind is the policy
% classification (clarification, local-correction, scope-restriction,
% unfaithful, false-source, open-gap); status is open, wip (elimination
% underway), resolved, or historical. Machine-read by the site index;
% prints a verdict line.
\newcommand{\gapnote}[2]{\par\noindent\textbf{Verdict:} #1 (#2).\par}


\title{Common Blocking and Span Equality in the Non-periodic MPS Fundamental Theorem}
\author{}
\date{2026-04-30}

\begin{document}
\maketitle

\section{The compressed source argument}

The proof of the non-periodic matrix product state fundamental theorem passes
from arbitrary tensors with the same matrix product vectors to
basis-of-normal-tensors comparison data. The source argument is the MPDO
canonical-form paper \cite{Cirac2016MPDO_arXiv,Cirac2017MPDO_AnnPhys}, together
with the later review \cite{Cirac2021Matrix}.\footnote{For traceability, this
note corresponds to \href{https://github.com/LionSR/TNLean/issues/1047}{issue \#1047}.}
The cited papers are not being questioned. The point is that a formal proof has
to separate several canonical-form steps which are compressed in the source
prose.

For a tensor $A=(A^i)_{i=1}^d$, the source paper defines the matrix product
vectors by
\[
  |V^{(N)}(A)\rangle
  = \sum_{i_1,\ldots,i_N=1}^d
      \operatorname{tr}(A^{i_1}\cdots A^{i_N})
      |i_1\rangle\otimes\cdots\otimes |i_N\rangle,
\]
for $N=1,2,3,\ldots$ \cite[definition of the matrix product vector family,
lines~130--138]{Cirac2016MPDO_arXiv}.
After removing invariant subspaces, the same paper writes the tensor in block form
\[
  A^i = \bigoplus_{k=1}^r \mu_k A_k^i,
\]
allowing zero blocks and scaling the nonzero blocks so that their transfer maps
have spectral radius one. Each nonzero irreducible block may still have peripheral
period. Blocking by a common multiple of the periods removes that periodicity
\cite[lines~214--231]{Cirac2016MPDO_arXiv}. The review follows the same
canonical-form argument \cite[lines~1798--1820]{Cirac2021Matrix}.

The later comparison is made through bases of normal tensors. A basis of normal
tensors is a minimal family $(A_j)_{j=1}^g$ of normal tensors such that every
matrix product vector of $A$ is a linear combination of the matrix product vectors
of the $A_j$, and such that those vectors are eventually linearly independent.
Equivalently, every normal block in the canonical form is a phase times a gauge
transform of exactly one basis tensor, with the family chosen minimally
\cite[Proposition~\texttt{prop:char-BNT}, lines~271--280]{Cirac2016MPDO_arXiv}.
The review states the same characterization
\cite[Proposition~\texttt{prop:char-BNT}, lines~1846--1859]{Cirac2021Matrix}.
For a tensor in canonical form one then has an expansion
\[
  |V^{(N)}(A)\rangle
    = \sum_{j=1}^g \left(\sum_{q=1}^{r_j} \mu_{j,q}^N\right)
        |V^{(N)}(A_j)\rangle.
\]
This is the basis-of-normal-tensors expansion used in the equal-case proof
\cite[lines~283--302]{Cirac2016MPDO_arXiv}. The review uses the same expansion
\cite[lines~1864--1885]{Cirac2021Matrix}.

Injectivity enters at a different point. The source first defines injective and
block-injective canonical form, and then invokes the quantum Wielandt theorem to
say that, after a bounded further blocking, every normal tensor becomes injective
\cite[lines~317--332]{Cirac2016MPDO_arXiv}. The review records the corresponding
$3D^5$ bound for block-injective canonical form
\cite[lines~2120--2124]{Cirac2021Matrix}.
This step is not the same as period removal: period removal gives primitive normal
blocks, while the overlap-rigidity arguments used later require one-site
injectivity, or an explicit replacement after further blocking.

Finally, the source comparison for the equal case is compact. Equality or
proportionality of the matrix product vector families first matches the two bases
of normal tensors by phase and gauge. After relabeling, the coefficient
multiplicities are compared through the power sums
\[
  \sum_{q=1}^{r_{a,j}} \mu_{j,q}^N
  = \sum_{q=1}^{r_{b,j}} (\nu_{j,q}e^{i\phi_j})^N,
  \qquad \text{for all } j \text{ and } N.
\]
The elementary power-sum lemma then recovers the multiplicities and the weights,
and the blockwise gauges are assembled into the global conjugating matrix
\cite[lines~1181--1192]{Cirac2016MPDO_arXiv}. The review states the resulting
proportional and equal MPV fundamental theorems in the same form
\cite[lines~1891--1900]{Cirac2021Matrix}.

\section{The formal argument}

The formal equality relation \leanid{SameMPV}\textsubscript{2} compares matrix product
vector coefficients for all finite words, including the empty word.
This convention is stronger than the positive-length family in the source
definition. An all-zero block of bond dimension $D_0$ contributes nothing at
positive length, but it contributes $D_0$ at length zero, because the empty product
has trace $D_0$. Consequently the formal proof separates the zero-tail
contribution from the nonzero part, proves positive-length equality for the
nonzero parts, and records the exact $N=0$ identity separately.\footnote{See
the zero-block separation section of the canonical-form chapter and the
corresponding zero-tail formalization in
\path{TNLean/MPS/CanonicalForm/SectorComparison/CommonSectorData.lean} and
\path{TNLean/MPS/CanonicalForm/SectorComparison/CommonSectorTransport.lean}.}

The common blocking step is also more explicit. For a single family of blocks,
the source phrase ``block by the least common multiple of the periods'' hides no
mathematical difficulty. In the two-tensor comparison, however, the formal proof
must choose one positive physical blocking length for both sides. For each
original block with period-removal length $m_k$, the common length is written as
$p=m_k e_k$. The tensor obtained by first blocking by $m_k$ and then by $e_k$
has physical labels given by nested words, while the tensor obtained by
blocking once by $p$ has physical labels in the alphabet of words of length $p$.
The formal statement therefore keeps the relabeling between these two alphabets as
part of the statement.\footnote{The common-sector data and the relabeling results
are recorded in \path{blueprint/src/chapter/ch09_canonical.tex:1645--1824} and
\path{blueprint/src/chapter/ch09_canonical.tex:3230--3284}. The main formal
statement is
\texttt{MPSTensor.fundamentalTheorem\_after\_blocking\_commonLength\_commonSector},
with the one-sided relabeling-compatible zero-tail statement
\texttt{MPSTensor.zeroTail\_commonFlat\_of\_reindexed}; see
\path{TNLean/MPS/CanonicalForm/SectorComparison/CommonSectorData.lean} and
\path{TNLean/MPS/CanonicalForm/SectorComparison/CommonSectorTransport.lean}.}
This is why the canonical-form theorem supplies relabeled common-sector
families and a conditional label-compatibility statement, rather than immediately
replacing every canonical nonzero part after blocking by its flattened sector tensor.

The formal comparison theorems downstream now use the coefficient-comparison route
for BNT sectors. The live data are the matched-sector phase relation, the
coefficient identities extracted from eventual linear independence, the
copy-weight matching obtained from the power sums, and the block-diagonal global
gauge. Any additional Wielandt blocking needed for one-site injectivity is a
separate refinement for arguments that require injective blocks; it is not part of
the current coefficient-comparison theorem.\footnote{The final comparison route is
recorded in \path{blueprint/src/chapter/ch11_fundamental_theorem_proof.tex:31--460}.}

The paper's compact BNT matching and power-sum comparison is no longer represented
by the removed comparison lemmas from the former span-based route. The
current formal comparison follows the source proof more directly: matched BNT
sectors give the phase relation
\[
  V^{(N)}(Q_k)_\sigma
    = \zeta_k^N V^{(N)}(P_{\beta(k)})_\sigma,
\]
eventual linear independence gives the coefficient identities
\[
  c_N^{(\beta(k))}(P)=\zeta_k^N c_N^{(k)}(Q),
\]
Newton--Girard power-sum comparison recovers the copy weights, and the matched
block gauges assemble into the direct-sum global gauge.\footnote{See
\path{blueprint/src/chapter/ch11_fundamental_theorem_proof.tex:31--460},
especially the sector coefficient identity, copy-weight matching, and assembled
global-gauge entries. The corresponding formal
objects live in
\path{TNLean/MPS/FundamentalTheorem/SectorBNT/CoeffIdentity.lean},
\path{TNLean/MPS/FundamentalTheorem/SectorBNT/CopyWeightMatching.lean},
\path{TNLean/MPS/FundamentalTheorem/SectorBNT/Fundamental.lean}, and
\path{TNLean/MPS/FundamentalTheorem/SectorBNT/FundamentalCoord.lean}.
The removed comparison lemmas from the former span-based route are no longer used
in the final comparison chapter; the remaining proof follows this
coefficient-comparison route.}
The coefficient-identity and copy-weight lemmas state precisely what comparison
data are needed; the proportional case still requires the coefficient comparison
for the phases supplied by the sector matching theorem.

\section{Why these distinctions belong together}

These distinctions are coupled in the single conversion from an arbitrary equality
of MPV families to the hypotheses of the BNT sector-comparison theorem. The
zero-tail convention determines what is known at length zero after each blocking.
The common blocking construction determines the physical alphabet and the relabeled
sector families to which the coefficient comparison applies. Any later injectivity
refinement is reserved for separate arguments, such as overlap-rigidity
statements that require one-site injective blocks; it is not a hypothesis of the
SectorBNT coefficient route. The matched-sector coefficient identities and
copy-weight data are constructed for the relabeled nonzero-sector families, not
for an earlier family before the common-alphabet relabeling.

For this reason common blocking and coefficient comparison should be kept in one
account. The formal proof needs the whole chain: remove the zero tail, choose one
common physical blocking length, account for the word-relabeling, keep later
overlap-rigidity requirements separate, and turn BNT matching data into the
sector coefficient identities and copy-weight equations used by the final gauge
comparison.
The blueprint states this chain as the connection between the canonical-form
chapter and the final comparison chapter.\footnote{See
\path{blueprint/src/chapter/ch11_fundamental_theorem_proof.tex:539--629} for the
current equal and proportional BNT canonical-form statements. The proportional
global gauge remains conditional on the coefficient identity recorded at
\path{blueprint/src/chapter/ch11_fundamental_theorem_proof.tex:361--388}; the
equal case uses
\path{blueprint/src/chapter/ch11_fundamental_theorem_proof.tex:503--537}.}

This note should also not be read as a claim that the same-side equal-norm scalar
obstruction discussed in the broader fundamental-theorem work is a flaw in the
CPSV or CPGSV arguments. That obstruction corrected an over-strong formal statement
which tried to identify an equal-norm class with a single basis tensor. The BNT
expansion in the papers already allows multiple basis tensors and multiplicities.
The formal argument therefore moved from naive norm-class grouping to phase-class
sector decompositions and common-cover data.

\section{Conclusion}

The conclusion is that the cited papers compress several standard canonical-form
operations; no counterexample is claimed. Zero or nilpotent contributions are
invisible to positive-length MPVs, periods are removed by blocking, injectivity is
obtained by a later Wielandt-type blocking, BNT bases are matched by phase and
gauge, and coefficient multiplicities are recovered from power sums. The formal
development exposes these operations as separate statements because the formal MPV
relation includes the empty word, the two sides must be placed over one explicitly
relabeled alphabet after blocking, and the downstream comparison theorem constructs
matched-sector coefficient identities and copy-weight equations after this
relabeling.

Thus the comparison gives neither a counterexample nor a theorem-level source
error. It explains why the formal proof cannot compress the argument into a
single invocation of the BNT uniqueness theorem. The relevant comparison points are
zero-tail cancellation, common primitive block construction, common blocking,
zero-tail transport, matched-sector coefficient identities, copy-weight matching,
and iterated-blocking relabeling.\footnote{For traceability, these are recorded in
\href{https://github.com/LionSR/TNLean/issues/652}{issue \#652},
\href{https://github.com/LionSR/TNLean/issues/942}{issue \#942},
\href{https://github.com/LionSR/TNLean/issues/969}{issue \#969},
\href{https://github.com/LionSR/TNLean/issues/970}{issue \#970},
\href{https://github.com/LionSR/TNLean/issues/971}{issue \#971},
\href{https://github.com/LionSR/TNLean/issues/944}{issue \#944}, and
\href{https://github.com/LionSR/TNLean/issues/990}{issue \#990}.}

\bibliographystyle{alpha}
\bibliography{references}

\end{document}
